%! ~~~ Packages Setup ~~~ 
\documentclass[]{../../../../tex/classes/styledArticle}
\usepackage{../../../../tex/packages/hebrewSupport}
\usepackage{../../../../tex/packages/mathShortcuts}
\usepackage{../../../../tex/packages/theoremsSupport}
\usepackage{../../../../tex/packages/enfancysections}

%! ~~~ Document ~~~

\author{שחר פרץ}
\title{\textit{לינארית 2א}}
\date{19 במאי 2025}
\begin{document}
    \maketitle
    \textbf{תזכורת: }המשלים הישר לתמ''ו ציקלי מקסימלי 
    
    {ציקלי מקסימלי: תהי $T \co V \to V$} ניל', ו־$U \subseteq V$ ציקלי מקסימלי אז קיים $W$ תמ''ו של $V$ תמ''ו $T$־אינ' כך ש־$T = U \oplus W$. 
    
    \cola{$T \co V \to V$ ט''ל ניל' אז $V$ אי־פריק ל־$T$ אמ''מ $V$ ציקלי. }
    \begin{proof}
        \begin{itemize}
            \item[$\implies$]הוכחנו בשיעור הקודם
            \item[$\impliedby$]נניח $V$ אי־פריק. אז קיים $U \subseteq V$ תמ''ו ציקלי מקסימלי. לפי המשפט קיים $W \subseteq V$ תמ''ו $T$־אינ' כך ש־$T = U \oplus W$. ידוע $U, W$ תמ''וים אינ'. אם $U = \{0\}$ אז $V = 0$ ובפרט ציקלי. אחרת, מאי־פריקות $V$ ל־$T$, נסיק ש־$W = \{0\}$ ולכן $V = U$ ציקלי. 
        \end{itemize}
    \end{proof}
    
    \theo{(משפט ג'ורדן למקרה של $T$ ניל') תהי $T \co V \to V$ ניל' אז קיים פירוק של $V$ לסכום ישר של $V = \bigoplus U_i$ כאשר $U_i$ הם $T$־ציקליים. }
    \begin{proof}
        נמצא ב־$V$ ציקלי מקסימילי. אז קיים $W \subseteq V$ תמ''ו $T$־שמור כך ש־$V = U_1 \oplus W$. ידוע $T_{|_W} \co W \to W$ ניל', וכעת באינדוקציה שלמה על $\dim V$. 
    \end{proof}
    
    גרסה שקולה למשפט ג'ורדן: \theo{עבוד $T \co V \to V$ ניל', קיים בסיס $B$ של $V$ שהוא איחור של שרשראות. }
    
    בסיס כזה נקרא בסיס מג'רדן ופירושו של דבר היא ש־: 
    \[ [T]_B = \pms{\square & 0 & \dots &0 \\ 0 & \square & \ddots & \vdots \\ \vdots & \ddots & \ddots & 0 \\ 0 & \cdots & 0 & \square} \]
    כאשר הבלוקים הם בלוקי ג'ורדן אלמנטרי. בלוק ג'ורדן מהצורה הבאה יקרא אלמנטרי ויסומן: 
    \[ J_n(0) = \pms{0 & 0 & \cdots & 0 \\ 1 & \ddots & \ddots & \vdots \\ \vdots & \ddots & 0 & 0 \\ 0 & 0 & 1 & 0} \]
    (יש ספרים בהם זה ה־transpose של זה). 
    
    \theo{עבור $T \co V \to V$ ניל', אז בכל הפירוקים של $V = \bigoplus U_i$ עבור $U_i$ ציקליים (אי־פריקים) אז מספר תתי־המרחב מממיד נתון הוא זהה עבור כל פירוק. }
    
    במילים אחרות: כל מטריצה מייצגת של ט''ל ניל' דומה למטריצת ג'ורדן יחידה, עד כדי שינוי סדר הבלוקים. 
    
    השקילות בין המשפט לבין ה''במילים אחרות'' נובע מזה שגודל בלוק פחות $1$ הוא הממד של התמ''ו שנפרש ע''י הוקטורים בעמודו תהללו. 
    
    \lem{נניח $V = \bigoplus_{i = 1}^{k} U_i$ כאשר $U_i$ הוא $T$־אינ' (לא נניח אפילו ניל'), אז: 
    \begin{enumerate}[A.]
        \item \hfil $\displaystyle T(V) = \bigoplus_{i = 1}^{k}T(U_i)$
        \item \hfil $\displaystyle \ker T = \bigoplus_{i = 1}^k \ker T(U_i)$
    \end{enumerate}}
    הוכחה: בש.ב. 
    
    במקרה הכללי, \defi{בלוק ג'ורדן אלנטרי עם ערך $\lg$ הוא בלוק מהצורה: 
        \[ J_n(\lg) = J_n(0) + \lg I_n = J_n(0) = \pms{\lg & 0 & \cdots & 0 \\ 1 & \ddots & \ddots & \vdots \\ \vdots & \ddots & \lg & 0 \\ 0 & 0 & 1 & \lg} \]
    }
    למה זה הגיוני? כי הסיבה שעשינו מתכתחילה רדוקציה ל־$T$ ניל' היא כי $T - \lg I$ היא $T$־אינוו' וניל', ועתה רק נותר להוסיף את ה־$\lg I$ חזרה לקבלת המקרה הכללי. 
    
    \begin{proof}
        באינדוקציה על $n  = n(T)$. 
        \begin{itemize}
            \item עבור $n = 1$, העתקת ה‏$0$. $V$ מתפרק לסכום ישר של מרחבים מממד $1$. 
            \item צעד, נניח נכונות עבור $n \in \N$. נניח ש־$n(T) = n + 1$. נסמן פירוק: 
            \[ V = \bigoplus_{i = 1}^{k} U_i = \bigoplus_{i = 1}^{\ml} W_i \]
            נסדר את $(u_i)_{i =1}^{k}$ לפי גודל מימד, ונניח שרשימת הגדלים היא: 
            \[ (\underbrace{1, 1, \dots 1 }_{\times s}< a_1 \le \dots \le a_{p}) \implies s + p = k \]
           רשימת הממדים מגודל $1$ ועוד כל השאר. נעשה כנ''ל עבור $(w_i)_{i = 1}^{\ml}$ ונקבל: 
            \[ (\underbrace{1, 1, \dots 1}_{\times t} < b_1 \le \dots \le b_r) \implies t + r = \ml \]
            ידוע:
            \[ T(v) = \bigoplus_{i = 1}^{k} T(U_i) = \bigoplus_{i = 1}^{k} T(W_i), \quad n\cl{T_{|_{T(v)}}} = n, \quad p = r, \quad \forall i \co a_i - 1 = b_i - 1 \implies a_i = b_i \]
            (הפירוק ל־$s$ ו־$t$ דרוש כדי שהפירוק לעיל לא יכלול אפסים כאשר מחילים את $T$)
            (ידוע $a_i - 1 = b_i - 1$ כי אינדקס הנילפוטנטיות קטן ב־1 בהחלת $T$)
            
            משפט הממדים השני אומר ש־: 
            \[ a_i = \dim U_i = \dim \ker T_{|_{U_i}} + \underbrace{\dim \Img T_{|_{U_i}}}_{a_i - 1} \implies \dim \ker T_{|_{U_i}} = 1 \]
            מהטענה השנייה בלמה: 
            \[ \begin{aligned}
                \ker T = \bigoplus_{i = 1}^{k} \ker T_{|_{U_i}} \implies \dim \ker T &= \sum_{i = 1}^k \dim \ker T_{|_{U_i}} = k \\
                &= \sum_{i = 1}^{k} \dim \ker T_{W_i} = \ml
            \end{aligned}\implies k = \ml \implies s = t \]
        \end{itemize}
    \end{proof}
    
    ``נראה לי שמי שסיכם את ההרצאה קצת חירטט את הסטודנטים ומי שסיכם את ההרצאה לא הבין את החרטוט'' – בן על הסיכום של הסטודנטים. 
    
    \ndoc
\end{document}
